\begin{definition}[3.2.5]
\label{IV.2.3.2.5-fr}
Sean $X$ un preesquema localmente noetheriano, $\mathcal F$ un
$\mathcal O_X$-módulo coherente. Se llama descomposición irredundante de
$\mathcal F$ a una familia $(\mathcal F_\alpha)_{\alpha\in I}$ de
$\mathcal O_X$-módulos cociente de $\mathcal F$ tal que los
$\mathcal F_\alpha$ sean irredundantes, que la familia
$(\operatorname{Supp}(\mathcal F_\alpha))$ sea localmente finita y que el
homomorfismo canónico
$\mathcal F\to\bigoplus_{\alpha\in I}\mathcal F_\alpha$ sea inyectivo. Se dice
que una descomposición de este tipo es reducida si los conjuntos
$\operatorname{Ass}(\mathcal F_\alpha)$ son distintos dos a dos y si no existe
ninguna parte $J\ne I$ tal que la subfamilia
$(\mathcal F_\alpha)_{\alpha\in J}$ sea una descomposición irredundante de
$\mathcal F$.
\end{definition}

Si $(\mathcal F_\alpha)_{\alpha\in I}$ es una descomposición irredundante
(resp. irredundante reducida) de $\mathcal F$ y si se pone
$\mathcal F_\alpha=\mathcal F/\mathcal G_\alpha$, se dice también que la
familia $(\mathcal G_\alpha)_{\alpha\in I}$ de sub-$\mathcal O_X$-módulos de
$\mathcal F$ es una \emph{descomposición primaria de $0$ en $\mathcal F$}~;
se observará que la hipótesis de inyectividad del homomorfismo canónico
$\mathcal F\to\bigoplus_{\alpha\in I}\mathcal F_\alpha$ equivale a la
condición $\bigcap_{\alpha\in I}\mathcal G_\alpha=0$.

Si $(\mathcal F_\alpha)$ es una descomposición irredundante de $\mathcal F$,
decir que es \emph{reducida} equivale a decir que los
$\operatorname{Ass}(\mathcal F_\alpha)$ son distintos dos a dos y están
contenidos en $\operatorname{Ass}(\mathcal F)$~; si
$\operatorname{Ass}(\mathcal F_\alpha)=\{x_\alpha\}$ para todo $\alpha\in I$,
$\alpha\leadsto x_\alpha$ es una \emph{biyección} de $I$ sobre
$\operatorname{Ass}(\mathcal F)$~: estas propiedades son, en efecto, locales y
resultan, por tanto, de Bourbaki, \emph{Alg. comm.}, cap.~IV, \S~2,
núm.~3, prop.~4.
